\section{Methods}

\subsection{Architecture}
\label{sec:arch}
We propose a novel method which encodes samples as classifications
mapped to contrastive vectors. This directly constrains a model's
representations. Moreover, classifications can natively be used for
steering.

The model (Figure~\ref{fig:onto-stack}) consists of $l$ $\dictenc$
layers writing into a shared residual accumulator, decoded by a single
linear map. Each layer (Figure~\ref{fig:dictenc}) is composed of a
classifier layer followed by a dictionary layer. 
The classifier splits the input into $h$ independent soft assignment
vectors, each over one head's $k$ entries.
The dictionary maps each entry to an atom.
Each layer writes a weighted average of these atoms to an accumulator.
At each layer, the accumulator $\mathbf{r}_\ell$ is decoded and the residual
$\mathbf{x} - \mathrm{dec}(\mathbf{r}_\ell)$ is passed to the next layer.
Deep supervision is used to constrain each layer to reconstructing the input.

Because entries within a head are mutually exclusive rather than independent,
each dictionary head serves as a contrastive set.
An auxiliary loss term minimizing the heads' support overlap $\omega$
(Eq.~\ref{eq:support}) pushes each head onto its own coordinates of the
dictionary space (Figure~\ref{fig:head-simplex}).
Because the dictionary returns a linear combination of representative vectors,
subspaces are convex.



%The mechanism is small enough to quote, excerpted from
%\texttt{dictblock.py} with dtype casts, intervention plumbing, and two
%live no-ops --- the identity output activation and the optional per-head
%scaling vector --- removed:
%
%\begin{lstlisting}[language=Python]
%def cluster(self, K, temperature=1.0):   # logits -> classifications,
%    return softmax(K / temperature)      # per head, over k entries
%
%def dicts(self):                         # abs() prevents subtractive
%    return jnp.abs(self.weights)         # ontofeatures; (h, k, d)
%
%def fwd(self, P):                        # convex combination per head,
%    return jnp.einsum("...hk,hkd->...d", # summed over heads
%                      P, self.dicts())
%
%def __call__(self, K, temperature=1.0):
%    return self.fwd(self.cluster(K, temperature))
%\end{lstlisting}

\begin{figure}[tp]
\centering
\resizebox{\linewidth}{!}{\begin{tikzpicture}[
    font=\footnotesize,
    >={Stealth[round,length=1.8mm,width=1.6mm]},
    blk/.style={draw, rounded corners=2pt, align=center,
                minimum height=8mm, inner xsep=4pt},
    hook/.style={draw, rounded corners=2pt, densely dashed, align=center,
                 inner sep=3pt, font=\scriptsize, black!55},
    dec/.style={draw, rounded corners=2pt, fill=black!8,
                minimum height=6.5mm, inner xsep=3pt},
    cond/.style={draw, rounded corners=2pt, fill=black!4, align=center,
                 inner sep=2.5pt, font=\scriptsize},
    add/.style={draw, circle, minimum size=4.5mm, inner sep=0pt},
    lab/.style={inner sep=1.5pt, font=\scriptsize},
    sig/.style={->, semithick},
    aux/.style={->, densely dashed, semithick, black!45},
  ]

  %% ---- layer row --------------------------------------------------
  \node[lab] (x)  at (-1.5,0) {$\mathbf{x}\in\mathbb{R}^{d}$};
  \node[blk] (D0) at ( 0.3,0) {$\mathtt{DictEnc}_0$};
  \node[blk] (D1) at ( 4.3,0) {$\mathtt{DictEnc}_1$};
  \node      (Dd) at ( 7.8,0) {$\cdots$};
  \node[blk] (DL) at ( 9.7,0) {$\mathtt{DictEnc}_{l-1}$};

  \draw[sig] (x) -- (D0);

  \node[hook] (noise) at (-1.5,2.6) {$+\,\varepsilon$};
  \draw[aux] (noise) -- (x);

  %% ---- residual accumulator (top rail, y = 2.6) --------------------
  \node[add] (A0) at ( 2.3,2.6) {$+$};
  \node[add] (A1) at ( 6.1,2.6) {$+$};
  \node[add] (AL) at (11.9,2.6) {$+$};

  \node[lab] (R0) at (0.6,2.6) {$\mathbf{r}_0\!=\!\mathbf{0}$};
  \draw[sig]      (R0) -- (A0);
  \draw[sig]      (A0) -- (A1);
  \draw[sig]      (A1) -- (AL);
  \draw[semithick] (AL) -- (12.35,2.6);

  %\node[lab] at (8.9,3.05) {$\mathbf{r}\in\mathbb{R}^{e}$: residual accumulator};

  %% ---- per-layer contributions g_l f'_l (enter from below) ----------
  \draw[sig] (D0.north) -- ++(0,1.1) -| (A0.south);
  \draw[sig] (D1.north) -- ++(0,1.1) -| (A1.south);
  \draw[sig] (DL.north) -- ++(0,1.1) -| (AL.south);
  \node[lab,anchor=west] at (0.38,1.05) {$g_0\mathbf{f}'_0$};
  \node[lab,anchor=west] at (4.38,1.05) {$g_1\mathbf{f}'_1$};
  \node[lab,anchor=west] at (9.78,1.05) {$g_{l-1}\mathbf{f}'_{l-1}$};

  %% ---- shared decoder, one decode per prefix -----------------------
  \node[dec] (Y1) at ( 2.75,-1.9) {$\mathrm{dec}$};
  \node[dec] (Y2) at ( 6.55,-1.9) {$\mathrm{dec}$};
  \node[dec] (YL) at (12.35,-1.9) {$\mathrm{dec}$};

  \draw[sig] ( 2.75,2.6) -- (Y1.north);
  \draw[sig] ( 6.55,2.6) -- (Y2.north);
  \draw[sig] (12.35,2.6) -- (YL.north);
  \fill ( 2.75,2.6) circle (1.1pt);
  \fill ( 6.55,2.6) circle (1.1pt);

  \node[lab,anchor=west] at ( 2.85,1.0) {$\mathbf{r}_1$};
  \node[lab,anchor=west] at ( 6.65,1.0) {$\mathbf{r}_2$};
  \node[lab,anchor=west] at (12.45,1.0) {$\mathbf{r}_{l}$};

  %% ---- deep supervision --------------------------------------------
  \node[lab] (L1) at ( 0.9,-1.9) {$\mathrm{MSE}(\mathbf{x},\hat{\mathbf{x}}_1)$};
  \node[lab] (L2) at ( 4.9,-1.9) {$\mathrm{MSE}(\mathbf{x},\hat{\mathbf{x}}_2)$};
  \node[lab] (LL) at (10.7,-1.9) {$\mathrm{MSE}(\mathbf{x},\hat{\mathbf{x}}_{l})$};
  \draw[aux] (Y1.west) -- (L1);
  \draw[aux] (Y2.west) -- (L2);
  \draw[aux] (YL.west) -- (LL);

  %% ---- residual forwarding ------------------------------------------
  \node[cond] (C1) at (2.75,-3.9)
  {$\mathrm{sg}\,(\mathbf{x}-\hat{\mathbf{x}}_1)$};
  \node[cond] (C2) at (6.55,-3.9)
  {$\mathrm{sg}\,(\mathbf{x}-\hat{\mathbf{x}}_2)$};
  \draw[sig] (Y1) -- node[right,lab]{$\hat{\mathbf{x}}_1$} (C1);
  \draw[sig] (Y2) -- node[right,lab]{$\hat{\mathbf{x}}_2$} (C2);

  \draw[sig] (C1.east) -- node[left,lab,pos=0.72]{$\mathbf{y}_1$} (D1.south);
  \draw[sig] (C2.east) -- node[left,lab,pos=0.72]{$\mathbf{y}_2$} (7.8,-0.5);

  %% ---- model output --------------------------------------------------
  \node[lab] (out) at (12.35,-3.9) {$\hat{\mathbf{x}}$};
  \draw[sig,thick] (YL) -- node[right,lab]{$\hat{\mathbf{x}}_{l}$} (out);

  %% ---- x broadcast to every conditioning step -------------------------
  \draw[aux] (x.south) -- (-1.5,-5.3) -- (2.75,-5.3) -- (C1.south);
  \draw[aux] (2.75,-5.3) -- (6.55,-5.3) -- (C2.south);
  \fill[black!45] (2.75,-5.3) circle (1.1pt);

\end{tikzpicture}

}
\caption[\texttt{Ontologizer} architecture]{\textbf{\texttt{Ontologizer} architecture.}
	An \texttt{Ontologizer} consists of $l$ stacked \texttt{DictEnc} layers.
	Each returns a feature vector $\mathbf{f}'_\ell\in\mathbb{R}^{e}$,
	which is written, scaled by its input's gain $g_\ell$, to a residual
	accumulator $\mathbf{r}$.
	For each layer $\ell = 0, \dots, l-1$,
	$\mathbf{r}_\ell = \mathbf{r}_{\ell-1} + g_\ell\,\mathbf{f}'_\ell$ with
	$\mathbf{r}_{-1} := \mathbf{0}$; a stop-gradient variant
	($\mathrm{sg}$ on $\mathbf{r}_{\ell-1}$, available but off in the live
	configuration) would prevent gradient communication between layers.
	$\mathrm{dec}$ is a linear decoder with shape $\mathbb{R}^{d \times e}$.
	It uses deep supervision to obtain a coarse-to-fine representation where
	the component encoded by previous layers is inaccessible to later layers.
	At each layer $\ell$, $\hat{\mathbf{x}}_\ell = \mathrm{dec}(\mathbf{r}_\ell)$.
	The same decoder is shared by all layers.
	The component left unexplained after layer $\ell$ is given by
	$\mathbf{y}_\ell = \mathbf{x} - \hat{\mathbf{x}}_\ell$.
	The input to layer $\ell+1$ is given by
	$\tilde{\mathbf{y}}_\ell
	  = \mathbf{y}_\ell / \mathrm{L2}(\mathbf{y}_\ell) \oplus [1]$
	where $\mathrm{L2}$ is the unit norm. This prevents exploding gradients from repeated
	application of bilinear layers, as each scales its input by $\lVert\cdot\rVert^{2}$.
	A constant coordinate $[1]$ is appended, giving a dimension of $d+1$.
	This prevents the sign information of $\mathbf{y}_\ell$ from being lost
	by the bilinear transformation.
	Effectively, it adds a bias to an otherwise unbiased linear form
	while preserving eigendecomposition.

	$\mathrm{MSE}(\mathbf{x},\mathrm{dec}(\mathbf{r}_\ell))$ is computed separately for each layer.

	The input-noise hook drawn on $\mathbf{x}$ is available but disabled
	in every model reported here.
}
\label{fig:onto-stack}
\end{figure}

\subsubsection{DictEnc}
\label{sec:dictenc}
Each $\dictenc$ consists of a multihead bilinear classifier
($\bilinearblock$) with weights
$[\mathsf{W}, \mathsf{V}] \in \mathbb{R}^{2 \times h \times k \times d}$
($d+1$ after the first layer, see Eq.~\ref{eq:y_norm})
and a multihead linear dictionary ($\dictblock$)
with weights $\mathsf{D} \in \mathbb{R}^{h \times k \times e}$.
It accepts a vector $\mathbf{x} \in \mathbb{R}^{d}$,
which may be a residual from a previous layer.
Algorithm~\ref{alg:dictenc} gives one layer's training pass, and
Algorithm~\ref{alg:ontologizer} the stack and its loss.

\begin{algorithm}[tp]
	\caption{$\dictenc$, one layer, in training (\texttt{DictEnc.withStats}).
	Without the router, $S$ is all ones and $\mathrm{L1}_S = 0$; without
	fibers there are no coordinates $\mathsf{C}$; without
	\texttt{resid\_gain} the input is not split and $g = 1$.}
	\label{alg:dictenc}
	\DontPrintSemicolon

	\KwIn{accumulator $R \in \mathbb{R}^{b \times e}$; layer input
		$E \in \mathbb{R}^{b \times d'}$, whose last $c$ coordinates are
		constant (\texttt{resid\_const})}
	\KwData{bilinear classifier (Eq.~\ref{eq:bilinearblock}); router
		(\texttt{scaled}); fiber readouts $A_i \in \mathbb{R}^{r \times d'}$;
		$\dictblock$ (Algorithm~\ref{alg:dictblock})}
	\Hyper{logit noise $\sigma_K$; fiber bound $\beta$; the
		$\dictblock$ hyperparameters}
	\KwOut{accumulator $R$; assignments
		$\mathsf{P} \in \mathbb{R}^{b \times h \times k}$; statistics
		row $\mathbf{m}$ (Eq.~\ref{eq:stats})}

	\Begin{
		$g_n \leftarrow \mathrm{sg}\big[\lVert E_{n,1:d'-c}\rVert\big]$;\quad
			$\mathbf{u}_n \leftarrow \big[E_{n,1:d'-c}/g_n;\,
			E_{n,d'-c+1:d'}\big]$
			\tcp*{gain and shape}
		$\mathsf{K}, \mathrm{L1}_K \leftarrow
			\mathrm{classifier}(\mathbf{u})$
			\tcp*{L1 of the gate factor}
		$\mathsf{K} \leftarrow \mathrm{noise}_K(\mathsf{K})$
			\tcp*{Eq.~\ref{eq:noise_normal} or~\ref{eq:noise_batchnorm}}
		\eIf{router}{
			$S, \mathrm{L1}_S \leftarrow \mathrm{router}(\mathbf{u})$;\quad
				$S \leftarrow \lvert S \rvert$
				\tcp*{unless \texttt{router\_signed}}
		}{
			$S \leftarrow \mathbf{1}$;\quad $\mathrm{L1}_S \leftarrow 0$\;
		}
		$\mathbf{c}_{ni} \leftarrow \beta \tanh(A_i \mathbf{u}_n / \beta)$
			\tcp*{fiber coordinates, Eq.~\ref{eq:fiber}}
		$F', \mathsf{P}, (\mathrm{L1}, H, \mathrm{cossim}_b,
			\mathbb{E}\kappa, \max\kappa, \mathrm{KL}_m, \sigma)
			\leftarrow \dictblock(\mathsf{K}, S, \mathsf{C})$\;
		$\mathrm{KL}_{\mathrm{pwak}}, \mathrm{L2}_{\mathrm{pwak}}
			\leftarrow \mathrm{pwak}(\mathsf{P}, \mathbf{u})$
			\tcp*{zero unless enabled}
		$R_n \leftarrow R_n + g_n\,\mathbf{f}'_n$
			\tcp*{the contribution, at the input's gain}
		\Return{$R$, $\mathsf{P}$,
			$\mathbf{m} = (\mathrm{L1}_K, \mathrm{L1}, H, \mathrm{cossim}_b,
			\mathbb{E}\kappa, \max\kappa, \mathrm{KL}_m,
			\mathrm{KL}_{\mathrm{pwak}}, \mathrm{L2}_{\mathrm{pwak}},
			\mathrm{L1}_S, \sigma)$}\;
	}
\end{algorithm}


\begin{figure}[tp]
\centering
\resizebox{\linewidth}{!}{\begin{tikzpicture}[
    font=\footnotesize,
    >={Stealth[round,length=1.8mm,width=1.6mm]},
    blk/.style={draw, rounded corners=2pt, align=center,
                minimum height=8mm, inner xsep=4pt},
    hook/.style={draw, rounded corners=2pt, densely dashed, align=center,
                 inner sep=3pt, font=\scriptsize, black!55},
    lab/.style={inner sep=1.5pt, font=\scriptsize, align=center},
    sig/.style={->, semithick, align=center},
    aux/.style={->, densely dashed, semithick, black!45},
  ]

  %% ---- main chain ----------------------------------------------------
  \node[lab] (E)  at (-0.6,0) {$\tilde{\mathbf{y}}$};
  \node[blk] (CL) at ( 1.2,0) {$\bilinearblock^{d \to h \times k}$};
  \node[blk] (SM) at ( 4.9,0) {$\mathrm{softmax}_k(\cdot/t)$};
  \node[blk] (DB) at ( 9.4,0) {$\dictblock^{h \times k \to h \times e}$};
  \node[blk] (SU) at (12.3,0) {$\sum^{h}_{i=1}$};
  \node[lab] (F)  at (13.6,0) {$F'$\\ pooled\\ features};

  \draw[sig] (E)  -- (CL);
  \draw[sig] (CL) -- node[below,lab]{$K$\\ logits}   (SM);
  \draw[sig] (SM) -- node[below,lab]{$P$\\ classifications}   (DB);
  \draw[sig] (DB) -- node[below,lab]{$F$\\ features} (SU);
  \draw[sig] (SU) -- (F);

  %% ---- training-time hooks and the intervention interface (above) ----
  \node[hook] (nk) at (3.0,1.5) {$+\,\varepsilon_K$,\ winner dropout};
  \node[hook] (iv) at (7.0,2.3) {\texttt{DictIntervention}\\
                                 set / add / sub / ablate / scale};
  \node[hook] (nf) at (11.4,1.5) {$+\,\varepsilon_F$};
  \draw[aux] (nk) -- (3.0,0.20);
  \draw[aux] (iv) -- (7.0,0.20);
  \draw[aux] (nf) -- (11.4,0.20);

  %% ---- where each regulariser attaches (below) -----------------------
  \node[hook] (lp) at (4.9,-1.8)
    {$H$\quad $\mathrm{KL}_m$\quad $\mathrm{cossim}_b$};
  \node[hook] (lf) at (9.4,-1.8)
    {$\mathrm{L1}$\quad $\sigma$};
  \draw[aux] (SM.south) -- (lp);
  \draw[aux] (DB.south) -- (lf);

\end{tikzpicture}

}
\caption[\texttt{DictEnc} architecture]{\textbf{\texttt{DictEnc} architecture.}
	Dashed lines above indicate perturbations to the data,
	including noising and causal interventions.
	Dashed lines below indicate statistics used as tuneable auxiliary loss functions.
	Each \texttt{DictEnc} consists of a multihead bilinear classifier
	(see §\ref{sec:bilinear})
	and a multihead linear dictionary (see §\ref{sec:dict}).
	It accepts a vector $\mathbf{x} \in \mathbb{R}^{d}$,
	which may be a residual from a previous layer.
	Optional noise terms $\varepsilon_K$ and $\varepsilon_F$ may be added
	(see §\ref{sec:noise}).
	Winner dropout sets each head's largest logit to $-\infty$
	with probability $p_{\mathrm{drop}}$.

	Soft assignments $P$ are given by $\mathrm{softmax}_{k}(K/t)$
	(see §\ref{sec:labels}).
	The architecture provides methods for causal interventions on assignments.
	From $P$ we collect per-sample entropy $H$ (Eq.~\ref{eq:entropy}),
	between-sample similarity $\mathrm{cossim}_b$ (Eq.~\ref{eq:cossim_b}),
	and batch mean entropy $\mathrm{KL}_m$ (Eq.~\ref{eq:KL}).

	A $\dictblock$ accepts assignments
	$P \in \mathbb{R}^{h \times k}$
	and returns features $F \in \mathbb{R}^{h \times e}$ (Eq.~\ref{eq:dictblock}).
	From $P$ and the atoms we collect the heads' support overlap $\omega$
	(Eq.~\ref{eq:support}), and from $F$ the sparsity penalty $\mathrm{L1}$
	(Eq.~\ref{eq:L1}).
	$F$ is summed over the head dimension to obtain
	pooled features $\mathbf{f}' \in \mathbb{R}^{e}$, which are added to
	the accumulator $\mathbf{r}$.
}
\label{fig:dictenc}
\end{figure}

\subsubsection{Classifier}
\label{sec:bilinear}
A $\bilinearblock$ accepts inputs $\mathbf{x} \in \mathbb{R}^d$.
It returns logits $K \in \mathbb{R}^{h \times k}$ given by

\begin{equation}
	\bilinearblock([\mathsf{W},\mathsf{V}],\mathbf{x}) :=
	\left(\sum^{d}_{q=1}W_{\cdot,\cdot,q}\,x_q\right)
	 \odot \left(\sum^{d}_{q=1}V_{\cdot,\cdot,q}\,x_q\right)
	 \label{eq:bilinearblock}
\end{equation}

An optional noise term $\varepsilon_K$ may be added.
Winner dropout sets each head's largest logit to $-\infty$
with probability $p_{\mathrm{drop}}$.

\subsubsection{Assignments}
\label{sec:labels}
Soft assignments $P$ are given by $\mathrm{softmax}_{k}(K/t)$
where $t$ is a temperature hyperparameter. Annealing it sharpens the
assignments; each model's schedule is in Appendix~\ref{app:configs}.
The architecture provides methods for causal interventions on assignments.
From $P$ we collect per-sample entropy $H$, in bits, so a head
saturates at $\log_2 k = 5$ bits,

\begin{equation}
	\mathrm{entropy}_k(\mathbf{p}) :=
	-\sum^{k}_{j=1}p_{j}\log_2 p_{j},
	\qquad
	H := \mathbb{E}_{n}\,\mathbb{E}_{i}\!\left[
	\mathrm{entropy}_k(\mathbf{p}_{ni})\right]
	\label{eq:entropy}
\end{equation}

between-sample similarity $\mathrm{cossim}_b$, the mean absolute cosine
similarity of the per-head features $\mathsf{F}$ over all ordered pairs
of samples,

\begin{equation}
	\mathrm{cossim}_b(\mathsf{P}) :=
	\mathbb{E}_{n,n'}\left[\left\lvert
	\frac{\langle \mathbf{f}_{n}, \mathbf{f}_{n'}\rangle}
	     {\lVert \mathbf{f}_{n}\rVert \lVert \mathbf{f}_{n'}\rVert}
	\right\rvert\right],
	\qquad
	\langle \mathbf{f}_{n}, \mathbf{f}_{n'}\rangle
	= \sum^{h}_{i=1}\mathbf{p}_{ni}^{\top}G_{i}\,\mathbf{p}_{n'i}
	\label{eq:cossim_b}
\end{equation}

where $\mathbf{f}_{ni} := \mathsf{F}_{ni\cdot} \in \mathbb{R}^{e}$ is
head $i$'s feature vector for sample $n$,
$\mathbf{f}_n = \map_{i=1}^{h}\mathbf{f}_{ni}$ is their concatenation,
and
$G_i = \lvert D_{i}\rvert \lvert D_{i}\rvert^{\top} \in \mathbb{R}^{k \times k}$
is head $i$'s dictionary Gram, with $D_i := \mathsf{D}_{i\cdot\cdot}$ its
$k \times e$ dictionary. The implementation evaluates
Equation~\ref{eq:cossim_b} in entry space through the $G_i$ rather than
materializing $\mathsf{F}$; the two are identical in value and gradient,
but $\mathcal{O}(hk^2e + b^2hk)$ instead of $\mathcal{O}(b^2he)$.

Finally, $\mathrm{KL}_m$ is the divergence of the batch-mean assignment
vector from uniform, in bits, averaged over heads.

\begin{equation}
	\mathrm{KL}_m(\mathsf{P}) := \mathbb{E}_{i}\!\left[\log_2 k -
	\mathrm{entropy}_k\!\left(\bar{\mathbf{p}}_i\right)\right],
	\qquad
	\bar{\mathbf{p}}_i := \mathbb{E}_{n}\,[\mathbf{p}_{ni}]
	\label{eq:KL}
\end{equation}

In practice $H$ is not needed, as including a
$\mathrm{cossim}_b$ loss also minimizes $H$.
$\mathrm{KL}_m$ serves to anchor the batch mean of each assignment vector
$\mathbf{p}_{ni}$ to a uniform distribution.
This results in a uniform $\mathbf{p}_{ni}$ code corresponding
to a maximally generic feature vector.

\subsubsection{Dictionary}
\label{sec:dict}
A $\dictblock$ accepts assignments
$P \in \mathbb{R}^{h \times k}$
and returns features $F \in \mathbb{R}^{h \times e}$, whose row
$\mathbf{f}_i$ is head $i$'s output, given by

\begin{equation}
	\dictblock(\mathsf{D},P)_{i\cdot} = \mathbf{f}_i :=
	\sum^{k}_{j=1} p_{ij}\,\mathbf{a}_{ij},
	\qquad
	\mathbf{a}_{ij} := \lvert \mathsf{D}_{ij\cdot}\rvert
	\label{eq:dictblock}
\end{equation}

where the atom $\mathbf{a}_{ij}$ is the absolute value of head $i$'s
dictionary row for entry $j$ ($\mathsf{D}_{ij\cdot}$ itself when the
dictionary is signed). Because atoms are absolute values
and $p_{ij} \in [0,1]$, features are strictly additive.
Algorithm~\ref{alg:dictblock} gives the full training pass, with the
regularizing perturbations and statistics around this lookup.

\begin{algorithm}[tp]
	\caption{$\dictblock$ in training (\texttt{DictBlock.withStats}). At
	inference the four stochastic steps (revival, winner dropout, head
	dropout, fiber dropout) and the feature noise are skipped. Every
	statistic is stop-gradiented unless its loss term is enabled.}
	\label{alg:dictblock}
	\DontPrintSemicolon

	\KwIn{logits $\mathsf{K} \in \mathbb{R}^{b \times h \times k}$;
		router scales $S \in \mathbb{R}_{\ge 0}^{b \times h}$ (all ones
		without a router); fiber coordinates
		$\mathsf{C} \in \mathbb{R}^{b \times h \times r}$ (none without
		fibers)}
	\KwData{dictionary $\mathsf{D} \in \mathbb{R}^{h \times k \times e}$;
		fiber bases $U_{ij} \in \mathbb{R}^{e \times r}$}
	\Hyper{temperature $t$; feature noise $\sigma_F$; winner dropout
		$p_{\mathrm{drop}}$; revival rate $p_{\mathrm{rev}}$ and threshold
		$\rho$; head dropout $p_{\mathrm{head}}$; fiber dropout
		$p_{\mathrm{fib}}$; entries kept per head $k_{\mathrm{sel}}$
		($1$ under \texttt{ste}/\texttt{argmax}, $m$ under
		\texttt{top}$m$, $k$ under softmax)}
	\KwOut{pooled features $F' \in \mathbb{R}^{b \times e}$; assignments
		$\mathsf{P} \in \mathbb{R}^{b \times h \times k}$; statistics
		$(\mathrm{L1}, H, \mathrm{cossim}_b, \mathbb{E}_i\kappa_i,
		\max_i \kappa_i, \mathrm{KL}_m, \sigma)$}

	\Begin{
		$\mathbf{d}_{ij} \leftarrow \lvert \mathsf{D}_{ij\cdot} \rvert$
			\tcp*{atoms; $\mathsf{D}_{ij\cdot}$ itself if signed}
		\If(\tcp*[f]{revive starved entries}){$k_{\mathrm{sel}} < k$}{
			$u_{ij} \leftarrow \mathbb{E}_n\,\mathbf{1}\{K_{nij}$ is among
				head $i$'s top $k_{\mathrm{sel}}$ logits$\}$\;
			\ForEach{$(n,i)$ with probability $p_{\mathrm{rev}}$}{
				draw $j$ uniformly from $\{j : u_{ij} <
					\rho\,k_{\mathrm{sel}}/k\}$, if any\;
				$K_{nij} \leftarrow \max_{j'} K_{nij'}$
					\tcp*{tie the leader}
			}
		}
		\ForEach(\tcp*[f]{winner dropout}){$(n,i)$ with probability
			$p_{\mathrm{drop}}$}{
			$K_{n,i,\arg\max_j K_{nij}} \leftarrow -\infty$\;
		}
		$\mathsf{P} \leftarrow \mathrm{select}_k(\mathsf{K}/t)$
			\tcp*{softmax, top-$m$ or straight-through}
		$H \leftarrow \mathbb{E}_{n,i}\,\mathrm{entropy}_k(\mathbf{p}_{ni})$
			\tcp*{Eq.~\ref{eq:entropy}}
		$\bar{\mathbf{p}}_i \leftarrow \mathrm{sg}[\mathbb{E}_n\,
			\mathbf{p}_{ni}]$ \tcp*{head origin}
		\ForEach(\tcp*[f]{head dropout}){$(n,i)$}{
			$\mathbf{p}'_{ni} \leftarrow \bar{\mathbf{p}}_i$ with
				probability $p_{\mathrm{head}}$, else
				$\bar{\mathbf{p}}_i + (\mathbf{p}_{ni} - \bar{\mathbf{p}}_i)
				/ (1 - p_{\mathrm{head}})$\;
		}
		\ForEach(\tcp*[f]{fiber dropout}){$(n,i)$ with probability
			$p_{\mathrm{fib}}$}{
			$\mathbf{c}_{ni} \leftarrow \mathbf{0}$\;
		}
		$\mathbf{f}_{ni} \leftarrow S_{ni} \sum_{j} p'_{nij}\,
			\mathbf{d}_{ij}$ \tcp*{Eq.~\ref{eq:dictblock}; interventions act
			on $\mathbf{p}'$}
		$\tilde{\mathsf{F}} \leftarrow \mathrm{noise}_{\mathrm{featvar}}
			(\mathsf{F})$ \tcp*{Eq.~\ref{eq:noise_featvar}}
		$\mathrm{cossim}_b \leftarrow \mathrm{cossim}_b(\mathsf{P})$
			\tcp*{Eq.~\ref{eq:cossim_b}, through the Grams}
		$\kappa_i \leftarrow \mathbb{E}_{j \ne j'}
			\cos(\mathbf{d}_{ij}, \mathbf{d}_{ij'})$
			\tcp*{within-head row cosine}
		$\mathrm{KL}_m \leftarrow \mathbb{E}_i \left[\log_2 k -
			\mathrm{entropy}_k(\bar{\mathbf{p}}_i)\right]$
			\tcp*{Eq.~\ref{eq:KL}}
		$\mathbf{s}_i \leftarrow \sum_j \bar{p}_{ij}\,
			\lvert\mathbf{d}_{ij}\rvert$;\quad
			$\sigma \leftarrow \mathbb{E}_{i \ne i'}
			\cos(\mathbf{s}_i, \mathbf{s}_{i'})$
			\tcp*{support overlap}
		$\mathrm{L1} \leftarrow \mathbb{E}_n \sum_i
			\lVert \tilde{\mathbf{f}}_{ni} \rVert_1$\;
		$\mathbf{f}'_n \leftarrow \sum_i \tilde{\mathbf{f}}_{ni}$
			\tcp*{concatenated, not summed, under $\mathtt{ConcatDictBlock}$}
		\If(\tcp*[f]{Eq.~\ref{eq:fiber}, then any norm cap}){fibers}{
			$\mathbf{f}'_n \leftarrow \mathbf{f}'_n + \sum_i S_{ni}
				\sum_j p'_{nij}\,U_{ij}\,\mathbf{c}_{ni}$\;
		}
		\Return{$F'$, $\mathsf{P}$,
			$(\mathrm{L1}, H, \mathrm{cossim}_b, \mathbb{E}_i\kappa_i,
			\max_i \kappa_i, \mathrm{KL}_m, \sigma)$}\;
	}
\end{algorithm}


From the assignments and the atoms we collect the heads' support overlap
$\omega$, the mean cosine over ordered pairs of distinct heads between
their usage-weighted coordinate profiles $\boldsymbol{\pi}_i$,

\begin{equation}
	\boldsymbol{\pi}_i := \sum^{k}_{j=1} \bar{p}_{ij}\,
	\lvert \mathbf{a}_{ij}\rvert,
	\qquad
	\omega := \mathbb{E}_{i \neq i'}\,
	\frac{\langle \boldsymbol{\pi}_i, \boldsymbol{\pi}_{i'}\rangle}
	     {\lVert \boldsymbol{\pi}_i\rVert \lVert \boldsymbol{\pi}_{i'}\rVert}
	\label{eq:support}
\end{equation}

with $\bar{\mathbf{p}}_i$ the batch-mean assignment of
Equation~\ref{eq:KL}. The overlap is zero exactly when no two heads share
a coordinate of the dictionary space and one when every head uses the
same coordinates in the same proportions (Section~\ref{sec:disjoint}). It
costs
$\mathcal{O}(hke + h^2e)$, independent of the batch. Weighting by usage
keeps dead entries out of it and gives the classifier a gradient path.
We also collect the sparsity penalty $\mathrm{L1}$ after the noise
$\varepsilon_F$ is added, $\tilde{\mathsf{F}} := \mathsf{F} +
\varepsilon_F$. $\mathrm{L1}$ is the batch mean of each sample's norm, so
it does not change with $b$ but scales with $h$ and $e$; its loss weight
is correspondingly small.

\begin{equation}
	\mathrm{L1}(\tilde{\mathsf{F}}) :=
	\frac{1}{b}\sum^{b}_{n=1}\sum^{h}_{i=1}\sum^{e}_{q=1}
	\lvert \tilde{f}_{niq}\rvert
	\label{eq:L1}
\end{equation}

$\mathrm{L1}$ tends to decrease with only $\mathrm{MSE}$ loss, but
incorporating it into the loss speeds up convergence.

\paragraph{Disjoint support}
\label{sec:disjoint}
Between heads, orthogonality of nonnegative features is not a rotation
constraint but a support-partition constraint. Every $\mathbf{f}_{ni}$
lies in the nonnegative orthant of
$\mathbb{R}^e$: the dictionary contributes $\lvert \mathsf{D} \rvert \geq 0$ and
the assignments $P \geq 0$. For nonnegative $\mathbf{v}, \mathbf{v}'$ the
inner product $\langle \mathbf{v},\mathbf{v}'\rangle = \sum_q v_q v'_q$ is
a sum of nonnegative terms, so it vanishes if and only if every term
does,

\begin{equation}
	\langle \mathbf{v},\mathbf{v}'\rangle = 0
	\iff
	\mathrm{supp}(\mathbf{v}) \cap \mathrm{supp}(\mathbf{v}') = \emptyset
	\label{eq:disjoint}
\end{equation}

Signed vectors admit no such conclusion: there orthogonality is
cancellation, and two vectors with identical support can be orthogonal.
The support overlap of Equation~\ref{eq:support} sidesteps this. It
compares the absolute values of the atoms, which have disjoint supports
exactly when they are orthogonal whatever the atoms' signs, so it
measures disjointness under a signed dictionary too; and it is computed
from the atoms and the batch-mean assignment rather than from noised
features, which leave the orthant and the argument with it.

The constraint binds the dictionary, not merely one sample's features.
Because the assignments are a softmax, $p_{nij} > 0$ for every entry, and expanding
Equation~\ref{eq:dictblock} gives

\begin{equation}
	\langle \mathbf{f}_{ni},\mathbf{f}_{ni'}\rangle =
	\sum^{k}_{j=1}\sum^{k}_{j'=1} p_{nij}\,p_{ni'j'}
	\left\langle \mathbf{a}_{ij}, \mathbf{a}_{i'j'} \right\rangle
	\label{eq:cross_head}
\end{equation}

a strictly positive combination of all $k^2$ cross-head atom overlaps.
A single sample on which heads $i$ and $i'$ are orthogonal therefore
forces $\langle \mathbf{a}_{ij}, \mathbf{a}_{i'j'}\rangle = 0$ for every
pair of atoms $(j,j')$: the two heads own disjoint blocks of coordinates
as a property of $\mathsf{D}$, holding for all inputs rather than on the
data that happened to be measured. That is the premise a static reading
of the dictionary needs, and it is stronger than what any signed-feature
architecture can ask for. The budget is not tight either. In the softmax
configuration (Table~\ref{tab:cfg-softmax}), $h{=}32$ heads partitioning
$e{=}2048$ coordinates leaves
64 dimensions per head to hold $k{=}32$ atoms, which need not be
disjoint from each other within a head; head-disjointness is most of
what the width $e > d$ buys.

Two different orthogonalities are in play and only the between-head one
is architectural. Disjoint support assigns each head a block of
coordinates; \emph{within} a head the $k$ atoms are contrastive
alternatives spread apart inside that block
(Figure~\ref{fig:head-simplex}). Section~\ref{sec:prelim-orth} measures
the second. Neither is guaranteed at convergence: $\omega$ is a penalty,
and the geometry is built to make a support partition the cheap solution
rather than to force one. Measured on the trained models, the cheap
solution was not taken: the within-head orthogonality holds, and no model
partitions its coordinates between heads (Appendix~\ref{sec:support}).
Under $\mathtt{ConcatDictBlock}$ heads write disjoint slices by
construction and $\omega$ is identically zero.

We add noise $\varepsilon_F$ to $F$, then sum over the head dimension to obtain
pooled features $\mathbf{f}' \in \mathbb{R}^{e}$, which are added to the
accumulator $\mathbf{r}$.

\begin{equation}
	\mathbf{f}' = \sum^{h}_{i=1}\tilde{\mathbf{f}}_{i}
	\label{eq:pool}
\end{equation}

\begin{figure}[tp]
\centering
\begin{tikzpicture}[
    font=\footnotesize,
    >={Stealth[round,length=1.8mm,width=1.6mm]},
    lab/.style={inner sep=1.5pt, font=\scriptsize},
  ]
  %% ---- head i --------------------------------------------------------
  \fill[black!8]   (0,0) -- (2.4,0) -- (1.2,2.0) -- cycle;
  \draw[semithick] (0,0) -- (2.4,0) -- (1.2,2.0) -- cycle;
  \fill (0,0)     circle (1.5pt) node[below,lab]{$\mathbf{a}_{i1}$};
  \fill (2.4,0)   circle (1.5pt) node[below,lab]{$\mathbf{a}_{i2}$};
  \fill (1.2,2.0) circle (1.5pt) node[above,lab]{$\mathbf{a}_{i3}$};
  \fill (1.05,0.72) circle (1.7pt) node[left,lab]{$\mathbf{f}_i$};

  %% ---- head i' ------------------------------------------------------
  \fill[black!8]   (5.0,0) -- (7.4,0) -- (6.2,2.0) -- cycle;
  \draw[semithick] (5.0,0) -- (7.4,0) -- (6.2,2.0) -- cycle;
  \fill (5.0,0)   circle (1.5pt) node[below,lab]{$\mathbf{a}_{i'1}$};
  \fill (7.4,0)   circle (1.5pt) node[below,lab]{$\mathbf{a}_{i'2}$};
  \fill (6.2,2.0) circle (1.5pt) node[above,lab]{$\mathbf{a}_{i'3}$};
  \fill (6.55,0.60) circle (1.7pt) node[right,lab]{$\mathbf{f}_{i'}$};

  \draw[<->, densely dashed, semithick, black!55] (1.35,0.71) -- (6.30,0.61);
  \node[lab] at (3.85,1.35)
    {$\omega$ penalizes heads $i$, $i'$ sharing coordinates};

  \node[lab, anchor=north, align=center] at (3.7,-0.9)
  {$\mathbf{f}'=\sum^h_{i=1} \mathbf{f}_i$,\qquad
  $\mathbf{f}_i=\sum^k_{j=1} p_{ij}\,\mathbf{a}_{ij}$,\qquad $\sum^k_{j=1} p_{ij}=1$};
\end{tikzpicture}


\caption[Per-head geometry]{\texttt{DictBlock} \textbf{head geometry} for $h{=}2$, $k{=}3$.
	The output space of a head $i$ forms a simplex with its atoms
	$\mathbf{a}_{ij}$, the rows of $\lvert D_i\rvert$, as the vertices.
	The atoms form the spanning set of the
	feature space. Any $\mathbf{f}_i$ must be a linear combination of them.
	As $F$ becomes sparser, these vectors become more approximately
	orthogonal, allowing them to be treated as a basis.
	The assignment code $P$ corresponds to soft assignments to mutually
	contrastive categories rather than independently firing features.
	Because elements of $P$, $p_{ij} \in [0,1]$,
	outputs must be within the convex hull of head $i$'s atoms.
	Because $\lvert \mathsf{D}\rvert\geq0$ and $P\geq0$, $F\geq0$;
	therefore $\cos(\mathbf{f}_i, \mathbf{f}_{i'}) \in [0,1]$, and driving the support
	overlap $\omega$ to zero asks the heads to partition the
	coordinates of $\mathbb{R}^e$ between them
	(Section~\ref{sec:disjoint}).
}
\label{fig:head-simplex}
\end{figure}

\subsubsection{Decoder}
\label{sec:decoder}
Each $\dictenc$ output $\mathbf{f}'_\ell$, scaled by its input's gain
$g_\ell$, is added to an accumulator $\mathbf{r}$.

\begin{equation}
	\mathbf{r}_\ell = \sum_{\ell'=0}^{\ell} g_{\ell'}\,\mathbf{f}'_{\ell'},
	\qquad
	g_\ell := \mathrm{sg}\!\left[\mathrm{L2}(\mathbf{y}_{\ell-1})\right]
\end{equation}

Under \texttt{resid\_gain} a layer classifies only the direction of its
input, so the gain carries the magnitude it cannot see; without it
$g_\ell = 1$. Layer 0's input is the encoded input itself,
$\mathbf{y}_{-1} := \mathbf{x}$, whose gain is identically 1 on unit-norm
SONAR embeddings.

A single-layer unbiased linear decoder $\mathrm{dec}$ projects
$\mathbf{r}_\ell$ back to the input space $\mathbb{R}^d$. Weights
$W_{\mathrm{dec}} \in \mathbb{R}^{d \times e}$
are shared for all $\mathtt{DictEnc}_\ell$.

\begin{equation}
	\hat{\mathbf{x}}_\ell = W_{\mathrm{dec}}\,\mathbf{r}_\ell
	\label{eq:dec}
\end{equation}

The residual $\mathbf{y}_\ell$ is given by $\mathbf{x} - \hat{\mathbf{x}}_\ell$, stop-gradiented so that
lower layers cannot write a communication code into the residual instead
of reducing it.
It is normalized by the unit norm $\mathrm{L2}$, with $\epsilon$ the
floating-point epsilon of the working dtype.

\begin{equation}
	\mathbf{y}_\ell = \mathrm{sg}\!\left[\mathbf{x} - \hat{\mathbf{x}}_\ell\right],
	\qquad
	\mathrm{L2}(\mathbf{y}) := \sqrt{\sum^{d}_{q=1} y_q^2}
	\label{eq:L2}
\end{equation}

\begin{equation}
	\tilde{\mathbf{y}}_\ell =
	\frac{\mathbf{y}_\ell}{\mathrm{L2}(\mathbf{y}_\ell) + \epsilon} \oplus [1]
	\label{eq:y_norm}
\end{equation}

A constant coordinate $[1]$ is appended, giving a dimension of $d+1$.
This prevents the sign information of $\mathbf{y}_\ell$ from being lost
by the bilinear transformation.
Effectively, it adds a bias to an otherwise unbiased linear form
while preserving eigendecomposition.

\subsubsection{Noise}
\label{sec:noise}
Noise is added to $X$, $\mathsf{K}$, and $\mathsf{F}$, controlled by
hyperparameters $\sigma_X$, $\sigma_K$, and $\sigma_F$. Each draws a
standard normal of the same shape as its target,

\begin{equation}
	\varepsilon \sim \mathcal{N}(0,1)
\end{equation}

and scales it three different ways. Classifier logits take unscaled
Gaussian noise,

\begin{equation}
	\mathrm{noise}_{\mathrm{normal}}(\mathsf{K}) :=
	\mathsf{K} + \sigma_K\varepsilon
	\label{eq:noise_normal}
\end{equation}

inputs take noise proportional to their own norm, so the scale is a
fraction of signal magnitude rather than an absolute quantity,

\begin{equation}
	\mathrm{noise}_{\mathrm{batchnorm}}(X) := X + \sigma_X\,\varepsilon
	\odot \frac{\mathrm{sg}\!\left[\mathrm{L2}(X)\right]}{\sqrt{d}}
	\label{eq:noise_batchnorm}
\end{equation}

and features take noise proportional to each feature's standard
deviation across the batch.

\begin{equation}
	\mathrm{noise}_{\mathrm{featvar}}(\mathsf{F}) := \mathsf{F} +
	\sigma_F\,\varepsilon \odot \mathrm{sg}\!\left[
	\sqrt{\mathrm{var}_n(\mathsf{F}) + \epsilon}\right]
	\label{eq:noise_featvar}
\end{equation}

Here $\mathrm{L2}(X)$ is the per-sample norm of
Equation~\ref{eq:L2}, broadcast over $d$. Both scales are
stop-gradiented: without $\mathrm{sg}$ the model could evade its own
noise by shrinking feature variance, and the $\epsilon$ inside the root
guards the $0/0$ gradient where a feature's batch variance underflows to
exactly zero. Because $\sigma_K$ enters the softmax as $\sigma_K/t$,
holding it constant under a falling temperature would grow effective
exploration rather than hold it, so models with an annealed temperature
anneal $\sigma_K$ with it. Each model's noise settings are in
Appendix~\ref{app:configs}.

\subsubsection{Router and fibers}
\label{sec:router-fiber}
Two optional branches read the same shaped input $\mathbf{u}$ as the
classifier and meet it in the dictionary lookup
(Figure~\ref{fig:dictenc-full}). Both are optional; the models that use
them are listed in Appendix~\ref{app:configs}.

The \emph{router} (\texttt{scaled}) is a second bilinear map, from
$\mathbf{u}$ to one scale per head, $\boldsymbol{\gamma} \in \mathbb{R}^h$,
made non-negative by an absolute value
(\texttt{router\_signed} drops it). Its scale $\gamma_i$ multiplies head
$i$'s whole output, so
the classifier decides \emph{what} a head says and the router \emph{how
much}; its own L1, $\mathrm{L1}_S$, prices the router alone, where
$\mathrm{L1}$ on the output can be paid down by either factor.

\emph{Fibers} (\texttt{fiber\_rank} $r$) give every entry a local basis
$U_{ij} \in \mathbb{R}^{e \times r}$ (in its head's slice under
$\mathtt{ConcatDictBlock}$), and each head reads $r$ coordinates
$\mathbf{c}_i = \beta\tanh(A_i\mathbf{u}/\beta)$ from the input, so
\begin{equation}
	\mathbf{f}_i = \gamma_i \sum^{k}_{j=1} p_{ij}\,\big(\mathbf{a}_{ij}
	      + U_{ij}\,\mathbf{c}_i\big).
	\label{eq:fiber}
\end{equation}
Under a hard code the selected entry is a discrete base point and the
fiber a continuous correction within the plane its basis spans
(Figure~\ref{fig:fiber}). Unconstrained, the correction can take over and
leave the base meaningless, so the base is kept honest by fiber dropout, a
norm cap ($\lVert U_{ij}\mathbf{c}_i\rVert \le \rho$ per head, or
summed over heads), a base-only deep-supervision term (\texttt{base\_aux},
which also decodes the fiber-free output), or training the fibers after
freezing the base. A radial fiber ($r{=}1$, no basis) only rescales the
selected entry. Statistics describe the base output. Head dropout
(\texttt{p\_head\_drop}) is the third addition here: in training it
replaces a head's classification by the batch mean, its origin, so heads
cannot co-adapt and ablating one to its origin stays in distribution.

\begin{figure}[tp]
\centering
\resizebox{\linewidth}{!}{% DictEnc with the router (`scaled`) and fibers (`fiber_rank`). Three
% branches read the shaped input u and meet in the dictionary lookup:
% the classifier picks entries (P), the router scales each head (gamma),
% and the fiber read gives each head local coordinates (C). Dashed boxes above
% an arrow act on it in training; dashed boxes below are statistics.
\begin{tikzpicture}[
    font=\footnotesize,
    >={Stealth[round,length=1.8mm,width=1.6mm]},
    blk/.style={draw, rounded corners=2pt, align=center,
                minimum height=8mm, inner xsep=4pt},
    hook/.style={draw, rounded corners=2pt, densely dashed, align=center,
                 inner sep=3pt, font=\scriptsize, black!55},
    mul/.style={draw, circle, minimum size=4.5mm, inner sep=0pt},
    lab/.style={inner sep=1.5pt, font=\scriptsize, align=center},
    sig/.style={->, semithick, align=center},
    aux/.style={->, densely dashed, semithick, black!45},
  ]

  %% ---- gain-shape split ------------------------------------------------
  \node[lab] (E)  at (-1.4,0) {$\mathbf{y}$};
  \node[blk] (GS) at ( 0.6,0) {gain--shape\\
                               $\mathbf{u}=\mathbf{y}/\lVert\mathbf{y}\rVert\oplus 1$};
  \draw[sig] (E) -- (GS);
  \coordinate (B) at (2.3,0);
  \draw[semithick] (GS.east) -- node[above,lab,pos=0.55]{$\mathbf{u}$} (B);
  \fill (B) circle (1.1pt);

  %% ---- three branches from u ---------------------------------------------
  % router: one non-negative scale per head
  \node[blk] (RT) at (4.4, 2.2) {$\bilinear^{d \to h}$\\ router};
  \node[blk] (AB) at (7.4, 2.2) {$\lvert\cdot\rvert$};
  % classifier: which entry, per head
  \node[blk] (CL) at (4.4, 0)   {$\bilinearblock^{d \to h \times k}$\\ classifier};
  \node[blk] (SM) at (7.4, 0)   {$\mathrm{softmax}_k(\cdot/t)$\\ or STE};
  % fiber read: r local coordinates per head
  \node[blk] (FR) at (4.4,-2.2) {$\mathsf{A} \in \mathbb{R}^{h \times r \times d}$\\ fiber read};
  \node[blk] (TB) at (7.4,-2.2) {$\beta\tanh(\cdot/\beta)$};

  \draw[sig] (B) |- (RT.west);
  \draw[sig] (B) -- (CL.west);
  \draw[sig] (B) |- (FR.west);
  \draw[sig] (RT) -- (AB);
  \draw[sig] (CL) -- node[below,lab]{$K$} (SM);
  \draw[sig] (FR) -- (TB);

  %% ---- the lookup ----------------------------------------------------------
  \node[blk, minimum height=5.6cm, text width=3.5cm] (DB) at (11.7,0)
    {$\dictblock$\\[6pt]
     $\mathbf{f}_i = \gamma_i \sum_{j} p_{ij}\,\big(\mathbf{a}_{ij}
            + U_{ij}\,\mathbf{c}_i\big)$\\[6pt]
     {\scriptsize atom $\mathbf{a}_{ij}$, fiber basis
      $U_{ij}\in\mathbb{R}^{e \times r}$}};

  \draw[sig] (AB) -- node[above,lab,pos=0.35]{$\boldsymbol{\gamma}\in\mathbb{R}^{h}$} (AB -| DB.west);
  \draw[sig] (SM) -- node[below,lab,pos=0.3]{$P\in\Delta_k^{h}$} (DB.west |- SM);
  \draw[sig] (TB) -- node[above,lab,pos=0.35]{$C\in\mathbb{R}^{h \times r}$} (TB -| DB.west);

  %% ---- combine heads, restore the gain --------------------------------------
  \node[blk] (CB) at (15.5,0) {$\sum_i$ \ or\ $\bigoplus_i$\\ {\scriptsize (concat)}};
  \node[mul] (MG) at (17.3,0) {$\times$};
  \node[lab] (F)  at (18.6,0) {$g\,\mathbf{f}'$\\ layer\\ output};
  \draw[sig] (DB.east |- CB) -- node[below,lab]{$F$} (CB);
  \draw[sig] (CB) -- node[above,lab]{$\mathbf{f}'$} (MG);
  \draw[sig] (MG) -- (F);
  \draw[sig] (GS.south) -- (0.6,-4.1) -| node[pos=0.25,below,lab]
       {$g=\lVert\mathbf{y}\rVert$, stop-gradient} (MG.south);

  %% ---- training-time hooks and the intervention interface (above) -------------
  \node[hook] (nk) at (5.9, 1.0) {$+\,\varepsilon_K$, winner dropout};
  \draw[aux] (nk) -- (5.9,0.2);
  \node[hook] (hd) at (8.95, 1.0) {head dropout\\ \texttt{DictIntervention}};
  \draw[aux] (hd) -- (8.95,0.2);
  \node[hook] (fd) at (8.95,-3.2) {fiber dropout};
  \draw[aux] (fd) -- (8.95,-2.4);
  \node[hook] (nf) at (15.5, 1.6) {$+\,\varepsilon_F$ on the base;\\
                                   fiber norm cap};
  \draw[aux] (nf) -- (14.05,0.2);

  %% ---- statistics (below) ------------------------------------------------------
  \node[hook] (sp) at (7.4,-1.1)
    {$H$\quad $\mathrm{KL}_m$\quad $\mathrm{cossim}_b$};
  \draw[aux] (SM.south) -- (sp);
  \node[hook] (ss) at (7.4, 3.45) {$\mathrm{L1}_S$};
  \draw[aux] (AB.north) -- (ss);
  \node[hook] (sf) at (11.7,-3.45)
    {$\mathrm{L1}$\quad $\omega$\quad $\kappa_i$\\ (base output and atoms)};
  \draw[aux] (DB.south) -- (sf);
  \node[hook] (ba) at (15.5,-1.9) {fiber-free $\mathbf{f}'_{\mathrm{base}}$\\
                                   (\texttt{base\_aux})};
  \draw[aux] (DB.east |- ba) -- (ba);

\end{tikzpicture}
}
\caption[\texttt{DictEnc} with router and fibers]{\textbf{\texttt{DictEnc}
	with router and fibers.} The gain--shape split hands the unit input
	$\mathbf{u}$ to three branches: the classifier (which entry, $P$),
	the router (how much, $\boldsymbol{\gamma}$, one non-negative scale per
	head) and the fiber read (where within the entry's fiber, $C$, $r$
	coordinates per head). They meet in the lookup
	(Equation~\ref{eq:fiber}); heads are summed, or concatenated under
	$\mathtt{ConcatDictBlock}$, and the result is scaled by the measured
	gain $g$. Dashed boxes above an arrow act on it during training;
	dashed boxes below are statistics, which describe the fiber-free
	base output. That output is also returned on its own for
	\texttt{base\_aux}.}
\label{fig:dictenc-full}
\end{figure}

\begin{figure}[tp]
\centering
% One head's geometry with a fiber, extending tikz/dict.tex: the head
% selects entry |d_{i2}|; its fiber is the plane U_{i2} spans at that entry,
% and the coordinates c_i place the output within it. The dashed circle is
% the norm cap (`fiber_eps`); the dotted ray is the radial variant, which
% only rescales the entry. Coordinates are precomputed (no calc library):
% d2 = (3.6, 3.2), basis columns u1 = (0.9, 0.25), u2 = (-0.2, 0.8).
\begin{tikzpicture}[
    font=\footnotesize,
    >={Stealth[round,length=1.8mm,width=1.6mm]},
    lab/.style={inner sep=1.5pt, font=\scriptsize},
  ]
  %% ---- origin and the head's simplex ----------------------------------
  \fill (0,0) circle (1.2pt) node[below left,lab]{$\mathbf{0}$};
  \fill[black!8]   (1.0,2.6) -- (3.6,3.2) -- (2.6,0.9) -- cycle;
  \draw[semithick] (1.0,2.6) -- (3.6,3.2) -- (2.6,0.9) -- cycle;
  \fill (1.0,2.6) circle (1.5pt) node[left,lab]{$\lvert\mathbf{d}_{i1}\rvert$};
  \fill (2.6,0.9) circle (1.5pt) node[below right,lab]{$\lvert\mathbf{d}_{i3}\rvert$};

  %% ---- radial variant: a per-entry gain along the entry's own ray ------
  \draw[densely dotted, black!55] (0,0) -- (4.5,4.0)
    node[above right,lab,align=left]{radial:\\ $(1+c_i)\,\lvert\mathbf{d}_{i2}\rvert$};

  %% ---- the fiber at the selected entry ---------------------------------
  \fill[black!18, opacity=0.8]
    (4.02,3.83) -- (4.26,2.87) -- (3.18,2.57) -- (2.94,3.53) -- cycle;
  \draw[black!60]
    (4.02,3.83) -- (4.26,2.87) -- (3.18,2.57) -- (2.94,3.53) -- cycle;
  \draw[densely dashed, black!55] (3.6,3.2) circle (0.75);
  \node[lab, black!55] at (4.75,2.25) {$\lVert U_{i2}\mathbf{c}_i\rVert \le \varepsilon$};

  \draw[->, semithick] (3.6,3.2) -- (4.14,3.35)
    node[right,lab]{$U_{i2}\mathbf{e}_1$};
  \draw[->, semithick] (3.6,3.2) -- (3.48,3.68)
    node[above left,lab]{$U_{i2}\mathbf{e}_2$};

  \fill (3.6,3.2) circle (1.5pt);
  \node[lab, anchor=east] at (2.9,3.25) {$\lvert\mathbf{d}_{i2}\rvert$};

  %% ---- the output: base point plus its fiber offset ---------------------
  \draw[->, thick] (3.6,3.2) -- (3.98,3.605);
  \fill (3.98,3.605) circle (1.8pt);
  \node[lab, anchor=south west] at (3.98,3.62)
    {$F_i = \lvert\mathbf{d}_{i2}\rvert + U_{i2}\mathbf{c}_i$};

  %% ---- legend ------------------------------------------------------------
  \node[lab, anchor=north west, align=left] at (-0.3,-0.45)
    {$\mathbf{c}_i = b\tanh(A_i\mathbf{u}/b)$ is read from the layer input;\\
     the selected entry is the discrete base point, the fiber the continuous\\
     correction. Fiber-free, the head emits $\lvert\mathbf{d}_{i2}\rvert$:
     the base-only output.};
\end{tikzpicture}

\caption[Fiber geometry]{\textbf{Fiber geometry} for one head, extending
	Figure~\ref{fig:head-simplex}. The head selects the entry with atom
	$\mathbf{a}_{i2}$; its fiber is the plane spanned by
	$U_{i2}$ at that point, and the coordinates $\mathbf{c}_i$ place the
	output within it. The dashed circle is the per-head norm cap; the
	dotted ray is the radial variant, which moves the output only along
	the atom's own direction.}
\label{fig:fiber}
\end{figure}

\subsection{Training}
\label{sec:training}

\subsubsection{Embedding model}
\label{sec:sonar}
We embed texts using SONAR\cite{sonar}, a multimodal and
language-agnostic sentence embedding model. SONAR embeddings have
several desirable properties:
\begin{itemize}
	\item They correspond to entire texts rather than tokens.
	\item They are bidirectional.
	\item Language is encoded independently of text.
\end{itemize}
This allows us to validate steering fidelity against a ground truth.
Large Concept Models\cite{lcmteam2024largeconceptmodelslanguage}
generate directly in SONAR space, which makes it a natural substrate
for steering experiments.

\subsubsection{Loss function}
\label{sec:loss}
Reconstruction fidelity is trained with $\mathrm{MSE}$, where
$\mathbf{w} \in \mathbb{R}^d$ is a fixed per-dimension weight vector.

\begin{equation}
	\mathrm{MSE}(X,\hat{X}) := \frac{1}{bd}\sum^{b}_{n=1}\sum^{d}_{q=1}
	w_q\left(\hat{x}_{nq} - x_{nq}\right)^2
	\label{eq:mse}
\end{equation}

The total objective adds the ghost-gradient reconstruction and the
auxiliary statistics, each weighted by a scalar. Statistics are
\emph{summed} over the $l$ layers rather than averaged.

\begin{equation}
	L := \mathrm{MSE}(X,\hat{X})
	+ s_g\,\mathrm{MSE}(\hat{X}_g,\, X - \hat{X})
	+ \sum^{l-1}_{\ell=0}\mathbf{s}\cdot\boldsymbol{\psi}_\ell
	\label{eq:loss_total}
\end{equation}

Each layer contributes the statistics (Algorithms~\ref{alg:dictblock}
and~\ref{alg:dictenc})

\begin{equation}
	\boldsymbol{\psi}_\ell = \big[\,
	\mathrm{L1}_K \;\; \mathrm{L1} \;\; H \;\; \mathrm{cossim}_b \;\;
	\mathbb{E}_i\kappa_i \;\; \max_i \kappa_i \;\;
	\mathrm{KL}_m \;\; \mathrm{KL}_{\mathrm{pwak}} \;\; \mathrm{L2}_{\mathrm{pwak}}
	\;\; \mathrm{L1}_S \;\; \omega
	\,\big]^{\top}_\ell
	\label{eq:stats}
\end{equation}

where $\mathrm{L1}_K$ is the L1 of the classifier's gate, $\kappa_i$ the
mean within-head cosine between head $i$'s atoms,
$\mathrm{KL}_{\mathrm{pwak}}$ and $\mathrm{L2}_{\mathrm{pwak}}$ the two
partition-consensus statistics, and $\mathrm{L1}_S$ the router's L1. The
max of $\kappa_i$ is logged for a setpoint controller and never itself
penalized. Appendix~\ref{app:stats} gives the implementation's row
layout.

Each model's weights are in Appendix~\ref{app:configs}. The ghost path
is disabled throughout, since $\hat{X}_g$ fires only on features that
are exactly zero across the whole batch, which bilinear classifiers with
softmax selection and absolute-valued dictionaries never produce.

\begin{algorithm}[tp]
	\caption{$\ontologizer$ in training (\texttt{Ontologizer.withStats})
	and the loss it feeds (\texttt{Hyperparams.loss}), for the
	residual forward (\texttt{forward="resid"}) every configuration here
	uses. Without an encoder, $\mathrm{enc}$ is the identity; with direct
	atoms, $\mathrm{dec}$ is.}
	\label{alg:ontologizer}
	\DontPrintSemicolon

	\KwIn{batch $X \in \mathbb{R}^{b \times d}$}
	\KwData{encoder $\mathrm{enc}$; layers $\dictenc_1, \dots, \dictenc_l$
		(Algorithm~\ref{alg:dictenc}); decoder $\mathrm{dec}$
		(Eq.~\ref{eq:dec}); MSE weights $\mathbf{w}$; loss weights
		$\mathbf{s}$, $s_g$}
	\Hyper{input noise $\sigma_X$; deep supervision; \texttt{deepsup\_sg};
		\texttt{resid\_norm}, \texttt{resid\_const}, \texttt{resid\_gain}}
	\KwOut{loss $L$; per-layer statistics $\mathbf{m}_1, \dots,
		\mathbf{m}_l$}

	\Begin{
		$E \leftarrow \mathrm{enc}\big(\mathrm{noise}_X(X)\big)$
			\tcp*{Eq.~\ref{eq:noise_batchnorm}}
		$R \leftarrow \mathbf{0} \in \mathbb{R}^{b \times e}$;\quad
			$E \leftarrow E \oplus [1]$ \tcp*{if \texttt{resid\_const}}
		\For{$i \leftarrow 1$ \KwTo $l$}{
			$R_{\mathrm{prev}} \leftarrow R$\;
			$R, \mathsf{P}_i, \mathbf{m}_i \leftarrow \dictenc_i(R, E)$\;
			\If(\tcp*[f]{a decode per prefix}){deep supervision}{
				$\hat{X}_i \leftarrow \mathrm{dec}\big(R - R_{\mathrm{prev}}
					+ \mathrm{sg}[R_{\mathrm{prev}}]\big)$
					\tcp*{$\mathrm{dec}(R)$ without \texttt{deepsup\_sg}}
			}
			\If{$i < l$}{
				$Y \leftarrow \mathrm{sg}\big[X - \mathrm{dec}(R)\big]$
					\tcp*{Eq.~\ref{eq:L2}}
				$E \leftarrow Y / (\mathrm{L2}(Y) + \epsilon)$
					\tcp*{Eq.~\ref{eq:y_norm}; raw under
					\texttt{resid\_gain}}
				$E \leftarrow E \oplus [1]$
					\tcp*{if \texttt{resid\_const}}
			}
		}
		\lIf{not deep supervision}{$\hat{X}_l \leftarrow \mathrm{dec}(R)$}
		$\mathrm{MSE} \leftarrow \mathbb{E}_i\,
			\mathrm{MSE}_{\mathbf{w}}(X, \hat{X}_i)$
			\tcp*{Eq.~\ref{eq:mse}; each prefix weighted $1/l$}
		$L \leftarrow \mathrm{MSE} + s_g\,\mathrm{MSE}_g
			+ \mathbf{s} \cdot \sum_{i=1}^{l} \mathbf{m}_i$
			\tcp*{Eq.~\ref{eq:loss_total}; $\mathrm{MSE}_g = 0$ without
			ghosts}
		\Return{$L$, $(\mathbf{m}_i)_{i=1}^{l}$}\;
	}
\end{algorithm}


Reconstruction is trained with deep supervision:
$\mathrm{MSE}(X, \mathrm{dec}(R_\ell))$ is computed separately at
every layer and averaged, so each prefix carries weight $1/l$ and each
layer must explain what previous layers left
unexplained. The MSE is whitened: $\mathbf{w}$ in
Equation~\ref{eq:mse} is the per-dimension inverse variance, normalized
to mean one, the same whitening the evaluation applies. The
implementation applies $\sqrt{\mathbf{w}}$ to $X$ and
$\mathrm{dec}(R_\ell)$ before the squared error, which is equivalent.
Four statistics are active as loss terms in the reported models
(Figure~\ref{fig:dictenc}; Appendix~\ref{app:configs}):
between-sample similarity $\mathrm{cossim}_b$, computed from $P$
through the dictionary Gram so it measures feature-space similarity of
the per-head reconstructions, and which in practice also minimizes
per-sample entropy $H$; the batch-mean term $\mathrm{KL}_m$, anchoring
each head's mean
assignment vector to a uniform distribution so that a uniform code
corresponds to a maximally generic feature; an $\mathrm{L1}$ penalty on
features; and a between-head disjointness term, the support overlap
$\omega$, which measures the extent to which heads use nonoverlapping
coordinates of the dictionary space. Most reported models trained before
$\omega$ existed and used its predecessor (Appendix~\ref{app:stats}).
$\mathrm{L1}$ tends to fall under MSE alone; including it in the loss
speeds up convergence. Classification sharpness is controlled by the
temperature. Winner dropout and the noise terms
$\varepsilon_K$ and $\varepsilon_F$ regularize classifications and
features respectively.

\subsection{Evaluation}
Evaluation runs against $\topk$ SAEs trained on identical cached SONAR
embeddings under an identical whitened reconstruction objective. The
battery is deliberately SAEBench-style\cite{karvonen2025}: a
reconstruction-versus-capacity Pareto; categorical-form and
semantic-coherence tests of head structure; cross-model feature
splitting and seed stability; downstream text fidelity through the
paired M2M100 decoder; sparse language probing; shrinkage
decomposition; steering fidelity; and blinded auto-interpretability
scoring. A structural-ablation ladder of SAE variants, each adding one
Ontologizer-style commitment at a time, attributes effects to specific
structure. AxBench\cite{wu2025axbenchsteeringllmssimple} found that
simple supervised baselines beat SAEs at steering; the steering
comparison here is built to face that result.


